\subsection{Writing a Quadratic Function in Vertex Form} 
We can also complete the square to write a quadratic function in vertex form. However, in this case we are rewriting an \makeblank[1in], so we can't do things to ``both sides.'' To do this, we will need a new simplifying move:
\begin{fact}
We can add and subtract the same number to an expression without changing its value. That, is, for any expression $X$ and real number $c$,
$$X=X+c-c$$
for all values of any variables in $X$.
\end{fact}
Writing vertex form is easier for quadratic functions whose leading coefficient is 1.
\begin{Example}Write $f(x)=x^2+12x-7$ in vertex form. Give the coordinates of the vertex of its graph.
\begin{lpic}{}{5}
\regtextleft{0.25}{-4}{Vertex:};
\end{lpic}
\end{Example}
We can still write vertex form when the leading coefficient is not 1, but the algebra requires a little care.
\begin{Example}
Write $g(x)=-2x^2+8x-5$ in vertex form. Give the coordinates of the vertex of its graph.
\begin{lpic}{}{9}
\foreach \y/\lbl in {1/Group the variable terms,
										2/Factor out the coefficient,
										3/Complete the square,
										5/Distribute,
										6/Combine the constants,
										8/Vertex}
	\regtextleft{0.25}{-\y}{\lbl:};
\end{lpic}
\end{Example}